\documentclass[a4paper,fleqn,review]{els-cas-dc}

\usepackage[authoryear,longnamesfirst]{natbib}
\usepackage{booktabs}
\usepackage{amsmath}
\usepackage{amssymb}
\usepackage{array}
\usepackage{graphicx}
\usepackage{url}

\newcommand{\dspo}{\textsc{DSPO}}
\newcommand{\dspoplus}{\textsc{DSPO\_PLUS}}

\newcommand{\ArtifactTable}[2]{%
  \IfFileExists{#1}{\input{#1}}{%
    \begin{table}[t]
      \centering
      \caption{#2}
      \begin{tabular}{@{}l@{}}
        \toprule
        Pending regenerated artifact \\
        \bottomrule
      \end{tabular}
    \end{table}
  }%
}

\begin{document}
\let\WriteBookmarks\relax
\def\floatpagepagefraction{1}
\def\textpagefraction{.001}

\shorttitle{Behavior-aware service menu optimization}
\shortauthors{J. Ana}

\title[mode=title]{Behavior-Aware Service Menu Optimization for Many-to-One Demand-Responsive Transit with Meeting Points and Time Windows}

\author[1]{Jin Ana}
\ead{jandxx@buaa.edu.cn}

\affiliation[1]{organization={MoE Key Laboratory of Complex System Analysis and Management Decision, School of Economics and Management, Beihang University},
  city={Beijing},
  postcode={100191},
  country={China}}

\begin{abstract}
This paper studies many-to-one demand-responsive transit in which the platform
offers each passenger a limited service menu consisting of meeting-point,
pickup-time-window, and price alternatives. The research objective is to
compare no-pricing, static-pricing, \dspo, and \dspoplus{} policies under a
shared multinomial-logit choice model with an outside option, paired replayed
requests, and explicit provenance gates. The manuscript is currently written as
a claim-gated TR-C draft: the method, experimental design, and artifact contract
are specified, while empirical dominance claims remain blocked until checkpoint
provenance and formal evidence gates pass. This separation is deliberate. It
prevents diagnostic, placeholder, or failed rows from being converted into
publication claims and makes the eventual ranking test reproducible.
\end{abstract}

\begin{keywords}
demand-responsive transit \sep meeting points \sep service menu optimization \sep multinomial logit \sep decision-focused learning \sep time windows
\end{keywords}

\maketitle

\section{Introduction}

Demand-responsive transit (DRT) systems must balance operational efficiency with
passenger acceptance. Meeting points can reduce routing cost and improve vehicle
productivity, but they also impose walking and timing burdens on passengers
\citep{stiglic2015benefits,fielbaum2021ondemand,cortenbach2024}. A platform
therefore cannot treat a meeting-point assignment as a purely operational
decision. It must decide which alternatives passengers see, how those
alternatives are priced, and how time-window feasibility is communicated before
the passenger chooses or opts out.

We study this interaction as a behavior-aware service-menu problem. Each
displayed alternative is a bundle containing a pickup location, a pickup
time-window treatment, and a price. The platform selects a limited menu from a
larger feasible candidate set; the passenger then chooses among displayed
bundles and an outside option according to a multinomial-logit (MNL) model
\citep{mcfadden1974conditional,ben1985discrete}. The central tension is that the
menu should steer demand toward operationally efficient meeting points without
manufacturing profit gains through unrealistic price manipulation or by hiding
the opt-out channel.

The paper is organized around three contributions. First, it formulates
many-to-one DRT menu design as an assortment-style optimization problem with
route-aware costs, time-window filters, prices, and passenger choice. Second, it
positions \dspo{} and \dspoplus{} as decision-focused menu-learning policies
whose differences must be visible in the objective, cost-learning, and behavior
gating contracts rather than asserted post hoc. Third, it defines a reproducible
experimental ladder on the RC benchmark: no pricing, static pricing, \dspo, and
\dspoplus{} are compared under paired replay only after checkpoint provenance,
row schema, and claim-readiness gates pass.

The present draft intentionally separates manuscript structure from final
claims. Current artifacts report a blocked claim state because pilot/formal
checkpoint provenance is unavailable. As a result, the results section describes
the evidence contract, current blockers, and planned ranking tests rather than
claiming that \dspoplus{} already dominates the baselines.

\section{Literature Review}

\subsection{DRT, Dial-a-Ride, and Meeting Points}

Classical dial-a-ride and pickup-and-delivery models provide the operational
foundation for shared DRT systems \citep{cordeau2007dial,parragh2008survey,
ho2018survey}. In these models, routing, capacity, and time-window constraints
determine whether requests can be served. Recent work on meeting-point DRT
extends this view by allowing passengers to walk to shared pickup locations,
which can reduce vehicle distance and improve pooling efficiency
\citep{stiglic2015benefits,zheng2019meeting,fielbaum2021ondemand}. Our focus is
not only where meeting points are located, but which meeting-point alternatives
are displayed to a passenger at the decision epoch.

\subsection{Service Menus and Assortment Optimization}

The displayed menu is closely related to assortment optimization under discrete
choice \citep{talluri2004revenue,rusmevichientong2010dynamic,desir2022capacitated}.
In classical assortment settings, products have known utilities and revenues.
In DRT, each product is an endogenous service bundle whose cost depends on the
current routing state and whose attractiveness depends on walking, in-vehicle
time, time-window reliability, and price. This makes the menu problem both an
assortment decision and an operational control problem.

\subsection{SPO and Decision-Focused Learning}

Decision-focused learning trains predictive models through the downstream
optimization loss rather than through prediction error alone
\citep{elmachtoub2022smart}. The \dspo{} family used here follows this logic:
cost predictions matter because they change menus, prices, and passenger
choices. The \dspoplus{} extension is treated as an enhanced decision-focused
policy with behavior-gated menu objectives, opt-out guardrails, service-quit
penalties, and robust time-window metadata. The manuscript requires the
\dspo{}/\dspoplus{} distinction to be generated as \texttt{\detokenize{method_family}}
metadata in experiment rows, not inferred manually from policy labels.

\section{Problem Formulation}

Consider a set of passenger requests $\mathcal{R}$ arriving over a planning
horizon. For each request $r$, the platform constructs a feasible candidate set
$\mathcal{B}(r)$. A service bundle is
\begin{equation}
  b = (m_b, [t_b^{\mathrm{lo}}, t_b^{\mathrm{hi}}], p_b),
\end{equation}
where $m_b$ is a home or meeting-point pickup location,
$[t_b^{\mathrm{lo}}, t_b^{\mathrm{hi}}]$ is the pickup time window, and $p_b$
is the displayed price. The platform chooses a displayed menu
$\mathcal{M}(r) \subseteq \mathcal{B}(r)$ subject to a menu-size limit and
feasibility checks.

Passenger utility for a displayed bundle is modeled as
\begin{equation}
  U_{rb} = \alpha_0 + \beta_p p_b + \beta_{\mathrm{ivt}}\widehat{\tau}^{\mathrm{ivt}}_{rb}
  + \beta_{\mathrm{walk}} d_{rb} + \beta_{\mathrm{tw}} q_{rb},
  \label{eq:utility}
\end{equation}
where $d_{rb}$ is walking distance, $\widehat{\tau}^{\mathrm{ivt}}_{rb}$ is
predicted in-vehicle time, and $q_{rb}$ summarizes pickup-time-window
compatibility. In the implementation, $q_{rb}$ is the semantic counterpart of
time deviation, ETA/window filtering, and \texttt{menu\_eta\_filter\_mode}
diagnostics, not a second independent time-window model. The coefficient signs
are interpreted semantically: price, in-vehicle time, walking distance, and
time-window deviation reduce utility, while the base and home-pickup terms
capture alternative-specific preference shifts. The outside option is
normalized to $U_{r0}=0$ in the paper formula; generated rows still record the
actual runtime \texttt{\detokenize{outside_option_util}} value, which defaults to 0 and may
be disabled in diagnostic settings. The choice probability is
\begin{equation}
  P_{rb}(\mathcal{M}) =
  \frac{\exp(U_{rb})}{\exp(U_{r0}) + \sum_{j \in \mathcal{M}(r)} \exp(U_{rj})}.
  \label{eq:mnl}
\end{equation}

The platform's expected menu value is the MNL-weighted net contribution of the
displayed bundles,
\begin{equation}
  \widehat{\Pi}(\mathcal{M}(r)) =
  \sum_{b \in \mathcal{M}(r)} P_{rb}(\mathcal{M})
  \left(p_b - \widehat{c}_{rb}^{\mathrm{route}} - \widehat{c}_{rb}^{\mathrm{risk}}\right).
  \label{eq:menu_objective}
\end{equation}
Capacity, service feasibility, and time-window feasibility constrain which
bundles enter $\mathcal{B}(r)$ and which bundles remain eligible for display.

\section{Methodology}

\subsection{\dspo{} Baseline}

The \dspo{} baseline combines learned cost prediction with menu optimization.
It estimates operational costs for feasible bundles, constructs a limited menu,
and evaluates expected menu value under the MNL choice model. In the experiment
contract, \dspo{} rows are identified by generated
\texttt{\detokenize{method_family=DSPO}} metadata and must share request traces, seeds,
pricing mode, routing settings, and checkpoint provenance with all baselines.

\subsection{\dspoplus{} Extension}

\dspoplus{} is defined as the enhanced decision-focused policy family only when
the generated rows carry explicit \texttt{\detokenize{method_family=DSPO_PLUS}} metadata.
Its minimum contract includes behavior-gated menu objectives, outside-option
and opt-out guardrails, service-quit penalties, and robust time-window metadata.
Attention-based variants remain diagnostic/V2 and do not satisfy the v1
\dspoplus{} contract. Any intended advantage must come from the generated
family contract and downstream validation, not from a price-only shortcut or a
manual relabeling of current mainline rows.

\subsection{Menu and Time-Window Extension}

The Work2 contribution is the interaction of menus and time windows. A menu is
the displayed subset of candidate meeting-point bundles. A time-window filter is
an eligibility rule that can remove bundles before menu scoring. The no-filter
case is a diagnostic endpoint: it disables ETA pruning but does not disable
routing, capacity, or service feasibility. Formal claims must distinguish robust
time-window handling from diagnostic no-filter exposure.

\section{Experimental Design}

\subsection{Dataset}

The primary dataset is the RC benchmark because it is the controlled setting
used by the current Work2 robust-menu pipeline. Yanjiao or additional external
instances are future extensions unless the RC formal ladder passes first.

\subsection{Baselines and Policies}

The required comparison ladder is no pricing, static pricing, \dspo, and
\dspoplus. The DSPO-family variants are further split into clip and wide
configurations when the manifest supports both. The target ranking
\begin{equation}
  \dspoplus > \dspo > \text{Static Pricing} > \text{No Pricing}
\end{equation}
is a validation criterion, not an assumption. If the ranking fails, the workflow
stops and reports the failure reason, minimal fix, and rerun command.

\subsection{Metrics}

The primary metric is realized net profit. Secondary metrics include served
rate, accepted-home count, opt-out rate, meeting-point adoption rate,
price-boundary diagnostics, checkpoint load status,
\texttt{\detokenize{outside_option_util}}, \texttt{\detokenize{method_family}}, and artifact
eligibility. \texttt{\detokenize{home_share}} is retained as a compatibility field with a
total-choice denominator; it is not an accepted-home-only rate unless that
denominator is stated explicitly. Opt-out, accepted home pickup, and accepted
meeting-point pickup remain separate outcome states in simulator counters,
normalized rows, artifact gates, and manuscript language.

\section{Results}

\subsection{Current Artifact Status}

Current Work2 robust-menu artifacts are not claim-ready. The claim guard reports
blocked status due to missing checkpoint provenance and formal evidence not yet
completed. Therefore, this draft reports artifact status and planned result
families rather than empirical dominance claims.

\ArtifactTable{../work2_coding/artifacts/work2_robust_menu/tables/provenance_status.tex}{Current provenance and claim-gate status.}

\subsection{Planned Ranking Validation}

The formal result table will be generated only after no-pricing, static-pricing,
\dspo, and \dspoplus{} rows are regenerated under paired replay with loaded
checkpoint provenance. Until then, any ranking language remains conditional.

\ArtifactTable{../work2_coding/artifacts/work2_robust_menu/tables/policy_summary.tex}{Diagnostic policy summary from the current artifact bundle.}

\section{Ablation Study}

The ablation section will be populated after the corresponding experiment gates
pass. The required ablations are: (i) removing time-window filtering or
time-window feasibility effects, (ii) removing menu expansion so that the policy
cannot benefit from a richer candidate set, and (iii) decomposing the
\dspo{}--\dspoplus{} gap into cost-learning, behavior-gating, and pricing/menu
channels. Ablation outputs must be generated from artifacts, not edited by hand.

\section{Conclusion}

This draft reframes Work2 as a TR-C paper on behavior-aware service-menu
optimization for many-to-one DRT. The contribution is the joint treatment of
meeting points, time windows, pricing, and passenger choice under a
claim-gated, paired-replay experiment design. The final empirical conclusion
depends on passing the RC evidence ladder. Until those gates pass, the
manuscript's defensible claim is that the framework, contracts, and artifact
guards are in place; policy dominance remains to be verified.

\section*{GSD Execution Report}

Milestone v1.1 starts from the completed V1 robust-menu pipeline and continues
phase numbering from Phase 6. The immediate manuscript action in this draft is
format and structure scaffolding only. Formal paper claims are blocked until
the code/experiment audit, model consistency repair, baseline validation,
\dspo{} runs, and \dspoplus{} runs have passed.

\section*{Reviewer-Style Risk Analysis}

\textbf{Novelty risk.} The strongest novelty is the combination of
behavior-aware service menus, meeting points, time windows, and decision-focused
policy learning. The risk is that reviewers may see parts of the framework as
known unless the \dspo{}/\dspoplus{} distinction and menu-time-window
interaction are made concrete.

\textbf{Modeling risk.} The MNL parameters and outside option must be framed as
simulator design or calibrated inputs unless external passenger-demand
estimation is added.

\textbf{Experiment risk.} The target ranking cannot be asserted before
checkpoint provenance, paired replay, and artifact gates pass. Missing
checkpoint status is currently a blocking issue.

\textbf{Acceptance probability.} At this draft stage the acceptance probability
is low to medium because the paper structure is in place but formal empirical
support is not yet claim-ready. It can improve only after the RC evidence ladder
passes and reviewer-facing risks are addressed directly.

\printcredits

\bibliographystyle{cas-model2-names}
\bibliography{references}

\end{document}
